\documentclass[12pt, letterpaper]{report}
\usepackage{amsthm,amssymb}
\usepackage{mathtext}
\usepackage[T1]{fontenc}
\usepackage[utf8]{inputenc}
\usepackage{pgfplots}
\usepgfplotslibrary{colormaps}
\usepackage{sectsty}
\usepackage[explicit]{titlesec}
\pgfplotsset{compat=1.9}
\usepackage[english,russian]{babel}
\makeatletter
\renewcommand{\@chapapp}{Пункт}
\makeatother
\chaptertitlefont{\Large}
\begin{document}
\begin{titlepage}
\begin{center}
\vspace*{1cm}
\textbf{МОСКОВСКИЙ ФИЗИКО-ТЕХНИЧЕСКИЙ ИНСТИТУТ (НАЦИОНАЛЬНЫЙ ИССЛЕДОВАТЕЛЬСКИЙ УНИВЕРСИТЕТ)}\\
\vspace{0.5cm} Физтех-школа Радиотехники и компьютерных технологий\\
\vspace{5cm} \LARGE{Лабораторная работа 2.3.3\\
Исследование функции}
\vfill
\large{\textbf{Рогов Анатолий Б01-406}} \\
\large \today
\vspace{0.8cm}
\end{center}
\end{titlepage}
\chapter{Исходная функция}
\hfil $f(x) = {\sin({{30}\cdot{{x}^{3}})}}+{{(\cos({{{20}\cdot{x}}+{1})})}^{3}}$\\
\begin{figure}[h]
\centering
\begin{tikzpicture}
\begin{axis} [
	legend pos = north west,
	xlabel = {$x$},
	ylabel = {$f(x)$},
	width = 300,
	height = 300,
	restrict y to domain=-2:2,
	restrict x to domain=-3:3,
	grid = major,
	enlargelimits=true,
]
\legend{
	$f^{(0)}$
}
\addplot[blue, samples=750]{(((sin(deg((((30))*((((x))^((3)))))))))+((((cos(deg((((((20))*((x))))+((1)))))))^((3)))))};
\end{axis}
\end{tikzpicture}
\caption{График функции}
\end{figure}
\section{Производная $f^{(1)}$ по переменной "$x$"}
\renewcommand{\thesubsection}{\arabic{subsection}}
\titleformat{\subsection}{\normalfont\bfseries}{}{0em}{\thesubsection\ #1 }
\subsection{По свойству 32 замечания 333}
\begin{center} $({x}^{3})^{\prime} = {3}\cdot{{x}^{2}}$ \end{center} 

\subsection{В греческой и античной математике знали и доказали тот факт, что}
\begin{center} $({30}\cdot{{x}^{3}})^{\prime} = {30}\cdot{{3}\cdot{{x}^{2}}}$ \end{center} 

\subsection{Следующее утверждение непосредственно следует из определения 17}
\begin{center} $(\sin({{30}\cdot{{x}^{3}})})^{\prime} = {\cos({{30}\cdot{{x}^{3}})}}\cdot{{30}\cdot{{3}\cdot{{x}^{2}}}}$ \end{center} 

\subsection{По свойству 32 замечания 333}
\begin{center} $({20}\cdot{x})^{\prime} = 20$ \end{center} 

\subsection{Подобно тому как было установлено выше}
\begin{center} $({{20}\cdot{x}}+{1})^{\prime} = 20$ \end{center} 

\subsection{Последовательным применением теорем 1488 и 228 получаем}
\begin{center} $((\cos({{{20}\cdot{x}}+{1})}))^{\prime} = {{(-1)}\cdot{\sin({{{20}\cdot{x}}+{1})}}}\cdot{20}$ \end{center} 

\subsection{Тогда выполняется}
\begin{center} $({(\cos({{{20}\cdot{x}}+{1})})}^{3})^{\prime} = {{3}\cdot{{(\cos({{{20}\cdot{x}}+{1})})}^{2}}}\cdot{{{(-1)}\cdot{\sin({{{20}\cdot{x}}+{1})}}}\cdot{20}}$ \end{center} 

\subsection{По следствию из теоремы А. Петровича получаем}
\begin{center} $({\sin({{30}\cdot{{x}^{3}})}}+{{(\cos({{{20}\cdot{x}}+{1})})}^{3}})^{\prime} = {{\cos({{30}\cdot{{x}^{3}})}}\cdot{{30}\cdot{{3}\cdot{{x}^{2}}}}}+{{{3}\cdot{{(\cos({{{20}\cdot{x}}+{1})})}^{2}}}\cdot{{{(-1)}\cdot{\sin({{{20}\cdot{x}}+{1})}}}\cdot{20}}}$ \end{center} 

\titleformat{\subsection}{\normalfont\bfseries}{}{0em}{\thesubsection \ #1\ }
\subsection{Результат}
\begin{center} $f^{(0)}(x) = {\sin({{30}\cdot{{x}^{3}})}}+{{(\cos({{{20}\cdot{x}}+{1})})}^{3}}$ \end{center} 
\begin{center} $f^{(1)}(x) = {{\cos({{30}\cdot{{x}^{3}})}}\cdot{{30}\cdot{{3}\cdot{{x}^{2}}}}}+{{{3}\cdot{{(\cos({{{20}\cdot{x}}+{1})})}^{2}}}\cdot{{{(-1)}\cdot{\sin({{{20}\cdot{x}}+{1})}}}\cdot{20}}}$ \end{center} 
\chapter{Разложение по формуле Тейлора до \\ $o((x - x_0)^{1})$ в точке $x_0 = 0$}
\begin{center} $f^{0}(x) = {\sin({{30}\cdot{{x}^{3}})}}+{{(\cos({{{20}\cdot{x}}+{1})})}^{3}}$ \end{center}
\begin{center} $f^{0}(0) = 0.157729$ \end{center}
\section{Производная $f^{(1)}$ по переменной "$x$"}
\renewcommand{\thesubsection}{\arabic{subsection}}
\titleformat{\subsection}{\normalfont\bfseries}{}{0em}{\thesubsection\ #1 }
\subsection{Пусть сказанное выше чистейшая правда, тогда}
\begin{center} $({x}^{3})^{\prime} = {3}\cdot{{x}^{2}}$ \end{center} 

\subsection{В силу доказанного факта можем заметить}
\begin{center} $({30}\cdot{{x}^{3}})^{\prime} = {30}\cdot{{3}\cdot{{x}^{2}}}$ \end{center} 

\subsection{Для данного утверждения рассуждения аналогичны}
\begin{center} $(\sin({{30}\cdot{{x}^{3}})})^{\prime} = {\cos({{30}\cdot{{x}^{3}})}}\cdot{{30}\cdot{{3}\cdot{{x}^{2}}}}$ \end{center} 

\subsection{Очевидно, что}
\begin{center} $({20}\cdot{x})^{\prime} = 20$ \end{center} 

\subsection{По свойству 32 замечания 333}
\begin{center} $({{20}\cdot{x}}+{1})^{\prime} = 20$ \end{center} 

\subsection{Последовательным применением теорем 1488 и 228 получаем}
\begin{center} $((\cos({{{20}\cdot{x}}+{1})}))^{\prime} = {{(-1)}\cdot{\sin({{{20}\cdot{x}}+{1})}}}\cdot{20}$ \end{center} 

\subsection{Из определения 5 следует тот факт, что}
\begin{center} $({(\cos({{{20}\cdot{x}}+{1})})}^{3})^{\prime} = {{3}\cdot{{(\cos({{{20}\cdot{x}}+{1})})}^{2}}}\cdot{{{(-1)}\cdot{\sin({{{20}\cdot{x}}+{1})}}}\cdot{20}}$ \end{center} 

\subsection{Из полученного следует}
\begin{center} $({\sin({{30}\cdot{{x}^{3}})}}+{{(\cos({{{20}\cdot{x}}+{1})})}^{3}})^{\prime} = {{\cos({{30}\cdot{{x}^{3}})}}\cdot{{30}\cdot{{3}\cdot{{x}^{2}}}}}+{{{3}\cdot{{(\cos({{{20}\cdot{x}}+{1})})}^{2}}}\cdot{{{(-1)}\cdot{\sin({{{20}\cdot{x}}+{1})}}}\cdot{20}}}$ \end{center} 

\titleformat{\subsection}{\normalfont\bfseries}{}{0em}{\thesubsection \ #1\ }
\subsection{Результат}
\begin{center} $f^{(0)}(x) = {\sin({{30}\cdot{{x}^{3}})}}+{{(\cos({{{20}\cdot{x}}+{1})})}^{3}}$ \end{center} 
\begin{center} $f^{(1)}(x) = {{\cos({{30}\cdot{{x}^{3}})}}\cdot{{30}\cdot{{3}\cdot{{x}^{2}}}}}+{{{3}\cdot{{(\cos({{{20}\cdot{x}}+{1})})}^{2}}}\cdot{{{(-1)}\cdot{\sin({{{20}\cdot{x}}+{1})}}}\cdot{20}}}$ \end{center} 
\begin{center} $f^{1}(0) = -14.7389$ \end{center}
\section{Ответ}
\begin{center} $f(x) = 0.157729 + \frac{-14.7389}{1!}\cdot(x - 0)^{1} + o((x - 0)^{1})$ \end{center}
\section{График членов Тейлора}
\pgfplotsset{
	colormap={bright}{
	rgb255=(220,102,85),
	rgb255=(240,66,204)
}}
\begin{figure}[h]
\centering
\begin{tikzpicture}
\begin{axis} [
	legend pos = north west,
	restrict y to domain=-2:2,
	restrict x to domain=-3:3,
	xlabel = {$x$},
	ylabel = {$f(x)$},
	width = 300,
	height = 300,
	grid = major,
	enlargelimits=true,
	colormap name=bright,
	cycle list={[of colormap]},
	every axis plot/.append style={mark=none,ultra thick}
]
\addplot+ [samples=750]{0.157729};
\addplot+ [samples=750]{0.157729 + (-14.7389/1) * (x - 0) ^ 1};
\legend{
	$o({x^{0}})$,
	$o({x^{1}})$
}
\end{axis}
\end{tikzpicture}
\caption{График членов разложения}
\end{figure}
\end{document}
